\documentclass[10pt,twocolumn]{article}
\usepackage[T1]{fontenc}
\usepackage[utf8]{inputenc}
\usepackage{lmodern}
\usepackage[margin=1.8cm,top=2.4cm,bottom=2cm,columnsep=0.6cm]{geometry}
\usepackage{fancyhdr}
\usepackage{titlesec}
\usepackage{booktabs}
\usepackage{amsmath,amssymb,amsfonts}
\usepackage{microtype}
\usepackage{xcolor}
\usepackage{mdframed}
\usepackage{parskip}
\usepackage{enumitem}
\usepackage{hyperref}

\definecolor{rulered}{RGB}{180,30,30}
\definecolor{boxgray}{RGB}{245,245,245}
\definecolor{keywordgray}{RGB}{100,100,100}

\pagestyle{fancy}
\fancyhf{}
\fancyhead[L]{\textsc{\small Claude Mag}}
\fancyhead[C]{\small\textit{Machine Learning Research Digest}}
\fancyhead[R]{\small 2026-07-13}
\fancyfoot[C]{\small\thepage}
\renewcommand{\headrulewidth}{0.8pt}

\titleformat{\section}{\normalsize\bfseries\color{rulered}}{}{0pt}{}%
  [\vspace{-4pt}\color{rulered}\rule{\columnwidth}{0.4pt}]
\titlespacing{\section}{0pt}{8pt}{4pt}

\hypersetup{colorlinks=true,urlcolor=rulered,linkcolor=black}

\begin{document}
\setlength{\parindent}{0pt}
\setlength{\parskip}{4pt}

{\small\color{keywordgray}\textbf{Keywords:}
  diffusion models \textbullet{}
  random matrix theory \textbullet{}
  generative models \textbullet{}
  statistical learning}\par\vspace{4pt}

{\LARGE\bfseries\raggedright
  The Seed Knows the Model}\par\vspace{4pt}

{\large\itshape\raggedright
  Random Matrix Theory Reveals Why Diffusion Models
  Trained on Different Data Produce the Same Images}\par\vspace{6pt}

{\small\textbf{By} Binxu Wang, Jacob A. Zavatone-Veth \& Cengiz Pehlevan
  (Harvard Kempner Institute; Harvard Society of Fellows; Harvard SEAS)}\par\vspace{2pt}

{\small\color{keywordgray}arXiv:2602.02908 \textbullet{}
  ICML 2026 Oral / Honorable Mention}
\par\vspace{2pt}\noindent{\color{rulered}\rule{\columnwidth}{0.6pt}}\vspace{6pt}

Train two diffusion models on non-overlapping halves of ImageNet. Give
both the same noise seed. Despite never sharing a single training image,
the models produce strikingly similar outputs. This reproducible
observation has prompted concern about memorization and dataset leakage
---but the true cause is far more mundane and, paradoxically, more
fundamental: shared Gaussian statistics across any large natural image
dataset already constrain most of what a finite-data denoiser can
generate. A random matrix theory framework makes this precise,
predicting cross-split consistency from first principles.

\begin{mdframed}[backgroundcolor=boxgray,linecolor=rulered,linewidth=1pt,
    innerleftmargin=6pt,innerrightmargin=6pt,
    innertopmargin=6pt,innerbottommargin=6pt]
\textbf{\small AT A GLANCE}\par\vspace{4pt}
\begin{description}[leftmargin=0pt,labelsep=4pt,itemsep=2pt,parsep=0pt,
    font=\small\bfseries]
  \item[Problem:] Diffusion models trained on disjoint data subsets
    generate nearly identical images from the same noise seed, raising
    memorization and reproducibility concerns.
  \item[Why it matters:] Understanding the source of this consistency
    is critical for auditing generative models for data privacy, copyright,
    and reproducibility.
  \item[Method:] Random matrix theory analysis of the linear diffusion model;
    exact formulas for the expected denoiser and its cross-split variance.
  \item[Result:] Consistency traces to shared Gaussian statistics across
    splits; the effective noise level is renormalized as $\sigma^2 \to \kappa(\sigma^2)$
    via a self-consistent RMT equation.
  \item[Evidence:] RMT predictions match empirical cross-split cosine
    similarity within 2--4\% across all tested dataset sizes and aspect ratios.
\end{description}
\end{mdframed}

\section{The Big Picture}

Diffusion models generate images by iteratively denoising a noise vector
$x_T \sim \mathcal{N}(0,I)$. Because the entire generation process is
deterministic given the initial seed $x_T$, two models starting from the
same seed produce the same image if and only if they have learned the
same score function. The observation that models trained on \emph{different}
data nonetheless produce similar images therefore implies that the learned
score functions agree --- which in turn implies that the differences in
training data have largely not been captured.

Prior work attributed this phenomenon to memorization (the models share
training examples through implicit data augmentation) or to architectural
inductive biases (convolutional symmetries, attention patterns). This paper
shows neither explanation is necessary: the consistency is predicted by
the population-level Gaussian statistics of the data, which are approximately
identical across any two large-enough random subsets of natural images.

\section{Why Existing Approaches Fall Short}

Memorization-based accounts predict that cross-split consistency should
decrease as the two data splits become more diverse. However, the
consistency persists even when the splits are drawn from entirely
different semantic categories of ImageNet, falsifying this explanation.

Architecture-based accounts predict that consistency should vanish
for architectures trained without weight sharing (e.g., plain MLPs).
Empirically, consistency is observed across all tested architectures,
inconsistent with this account.

Statistical privacy audits (membership inference, differential privacy
bounds) focus on individual example memorization rather than
population-level structural similarities. None of these frameworks
could predict the direction or magnitude of cross-split agreement.

\section{Core Mechanism}

The consistency mechanism operates through the spectral structure of
the data covariance. In the linear diffusion model, the optimal denoiser
at noise level $\sigma$ is the Wiener filter:
\[
s_\sigma(x_\sigma) = -(\Sigma + \sigma^2 I)^{-1} x_\sigma
\]
where $\Sigma$ is the data covariance. Two models trained on different
splits both estimate $\Sigma$ from their respective finite datasets.
Their estimates $\hat{\Sigma}^{(1)}$ and $\hat{\Sigma}^{(2)}$ are noisy,
but their \emph{mean} over random splits converges to the \emph{same}
population covariance $\Sigma$. The shared population covariance is the
source of cross-split consistency: the two denoisers agree in expectation
and differ only in their finite-sample fluctuations, which vanish at
rate $O(1/\sqrt{N})$.

\section{Technical Formulation}

In the proportional limit $N, d \to \infty$ with $\gamma = d/N$ fixed,
the sample covariance resolvent $(\hat{\Sigma} + \sigma^2 I)^{-1}$ converges
to a deterministic function of the population covariance $\Sigma$ via
the self-consistent equation:
\[
\kappa(\sigma^2) = \sigma^2 + \gamma \int \frac{\lambda}{\lambda + \kappa(\sigma^2)}\,d\mu(\lambda)
\]
where $\mu$ is the spectral measure of $\Sigma$. This $\kappa(\sigma^2) > \sigma^2$
is the \emph{effective noise level}: finite-data denoisers experience more
noise than they receive, explaining the systematic bias toward the dataset
mean.

The expected cross-split agreement of generated samples is:
\[
\mathbb{E}[\langle x^{(1)}, x^{(2)} \rangle]
= \operatorname{tr}\!\left[\left(\Sigma + \kappa I\right)^{-2} \Sigma^2\right] / d
\]
Cross-split variance decomposes into three additive terms:
\[
\operatorname{Var}[\hat{s}_\sigma(x)] = V_{\text{aniso}} + V_{\text{inhom}} + V_{\text{size}}
\]
where $V_{\text{aniso}}$ captures eigenmode anisotropy, $V_{\text{inhom}}$
captures per-input inhomogeneity, and $V_{\text{size}} = O(1/N)$ is the
global $1/N$ scaling.

\section{Learning or Inference Procedure}

\begin{enumerate}[itemsep=2pt,parsep=0pt]
  \item Draw two independent data splits of size $N$ each from $p(x)$.
  \item Train a linear diffusion model (Wiener denoiser) on each split,
    obtaining score functions $\hat{s}_\sigma^{(1)}$ and $\hat{s}_\sigma^{(2)}$.
  \item Generate images $x^{(1)}$ and $x^{(2)}$ from the same noise seed
    $x_T$ using Euler--Maruyama sampling.
  \item Measure cross-split cosine similarity $\langle x^{(1)}, x^{(2)} \rangle
    / (\|x^{(1)}\|\|x^{(2)}\|)$.
  \item Compare to the RMT prediction derived from the population covariance
    $\Sigma$ and $\kappa(\sigma^2)$.
\end{enumerate}

The RMT prediction requires only $\Sigma$ and $\gamma = d/N$ --- neither
the specific split realizations nor any model training is needed.

\section{What the Guarantee Says}

The main theorem establishes: as $N, d \to \infty$ with $\gamma = d/N$
and signal-to-noise ratio bounded away from 0 and $\infty$:

(1)~\emph{Expectation:} $\mathbb{E}[\hat{s}_\sigma(x)]$ converges to
$-(\Sigma_{\text{eff}} + \kappa I)^{-1} x$ where $\Sigma_{\text{eff}}$
is a deterministic function of the population covariance.

(2)~\emph{Fluctuations:} $\operatorname{Var}[\hat{s}_\sigma(x)] = O(1/N)$,
so cross-split agreement is $1 - O(1/N)$ for fixed $d/N$.

(3)~\emph{Bias:} Generated samples are systematically pulled toward
the dataset mean along low-variance principal components, with magnitude
$O(\gamma)$ for fixed spectral gap.

\section{Experimental Findings}

\begin{center}
\small
\begin{tabular}{lcc}
\toprule
\textbf{Setting} & \textbf{Measured} & \textbf{RMT pred.} \\
\midrule
Linear, $N = 500$   & 0.47 & 0.46 \\
Linear, $N = 2000$  & 0.23 & 0.22 \\
$\gamma = 2.0$       & 0.61 & 0.60 \\
ImageNet PCA approx. & 0.38 & 0.36 \\
\bottomrule
\end{tabular}
\end{center}

All values are cross-split cosine similarity. RMT predictions match
observations within 2--4\% across all settings. The framework also
correctly predicts \emph{directional} bias: generated samples shift
toward the dataset mean along the eigenvectors with smallest
eigenvalues, which the Wiener filter systematically undershoots.

\section{Ablations and Interpretation}

\textbf{Dataset size:} Cross-split consistency decreases as
$O(1/\sqrt{N})$, confirming the $1/N$ fluctuation bound. Practitioners
can use this formula to determine the minimum dataset size needed to
keep consistency below an acceptable threshold.

\textbf{Aspect ratio $\gamma$:} Higher $\gamma$ (more parameters per
sample) increases cross-split consistency. Modern large-latent-space
models effectively operate at high $\gamma$, explaining why they show
stronger consistency than smaller models.

\textbf{Nonlinear diffusion:} Experiments on DDPM with U-Net score
models show qualitatively consistent results: cross-split similarity
decreases with $N$ and increases with $\gamma$, suggesting the linear
theory captures the dominant statistical effect even in nonlinear models.

\textbf{Semantic consistency:} Cross-split consistency in semantic
embedding space (CLIP) is 2--3$\times$ larger than in pixel space,
suggesting the linear mechanism dominates more strongly at the semantic
level where principal components carry more variance.

\section*{Reference}

Binxu Wang, Jacob A.\ Zavatone-Veth, Cengiz Pehlevan.
\textbf{A Random Matrix Theory Perspective on the Consistency of
Diffusion Models}.
arXiv:2602.02908, February 2026. ICML 2026 Oral / Outstanding Paper Honorable Mention.
\url{https://arxiv.org/abs/2602.02908}

\end{document}
